\chapter{Methodology and Data Collection}
\label{ch:method}

This chapter describes how library dependencies are extracted from
certification data and how the research questions are answered. The work
proceeds in four phases: dataset exploration and definition of the search
heuristics; the extraction pipeline with AI-assisted parsing; exploratory
analysis, metrics, and visualization; and packaging of the reproducible
notebooks.

\section{Data source and scope}
\label{sec:data}

The study covers all CC and FIPS~140-2 and 140-3 certificates indexed by
sec-certs.org. The processed datasets are loaded through the documented Python
interface~\cite{seccertsDocs}. The snapshot is downloaded once, with its
documents, and the notebooks then load it from local JSON files, as shown in
Listing~\ref{lst:datasets}.

\begin{lstlisting}[language=Python, caption={Loading the processed datasets},
  label={lst:datasets}]
from sec_certs.dataset.cc import CCDataset
from sec_certs.dataset.fips import FIPSDataset

from pathlib import Path

cc = CCDataset.from_json("./datasets/CC_dataset.json")
cc.root_dir = Path("./datasets")
fips = FIPSDataset.from_json("./datasets/FIPS_dataset.json")
fips.root_dir = Path("./datasets")
\end{lstlisting}

For each certificate, the pipeline uses the vendor, the product name, the
scheme, the certification dates, the certified security claims, and the text of
the associated documents: Security Targets, certification reports, and
maintenance reports for CC, and Security Policies for FIPS~140.

{Record the sec-certs version and the dataset snapshot date. Confirm the
exact certificate fields during Phase~1, because the API documentation does not
list the per-certificate FIPS fields.}

Because sec-certs.org is continuously updated, all counts in this thesis refer
to one pinned snapshot, which is archived with the toolkit.

\section{Detection heuristics}
\label{sec:heuristics}

The first extraction step is deliberately high in recall: a set of regular
expressions finds candidate mentions of each library and, where present, an
adjacent version number. PDF text extraction can split names across lines, so the
patterns tolerate whitespace, and they match version strings used across
OpenSSL's naming schemes, such as \texttt{0.9.8zh}, \texttt{1.0.2k}, and
\texttt{3.0.8}. Listing~\ref{lst:patterns} shows the initial patterns. BoringCrypto
is mapped to BoringSSL because it is the name of BoringSSL's FIPS module.
{Add a citation for BoringCrypto.}

\begin{lstlisting}[language=Python, caption={Initial library detection patterns},
  label={lst:patterns}]
import re

LIBRARIES = {
    "openssl": r"open\s?ssl",
    "boringssl": r"boring\s?(?:ssl|crypto)",
    "libressl": r"libre\s?ssl",
    "aws-lc": r"aws[\s-]?lc",
    "quictls": r"quictls",
}
QUALIFIER = (
    r"(?:[\s-]+(?:fips|object|module|"
    r"provider|library))*"
)
VERSION = (
    r"(?:[\s-]+(?:version\s+|v)?"
    r"(\d+\.\d+(?:\.\d+)?[a-z]{0,2}))?"
)
PATTERNS = {
    name: re.compile(
        rf"\b{lib}\b{QUALIFIER}{VERSION}",
        re.IGNORECASE,
    )
    for name, lib in LIBRARIES.items()
}
\end{lstlisting}

These patterns are a starting point. They will be tuned iteratively against the
data, especially to distinguish standard OpenSSL from vendor-branded wrappers and
independent forks. A match is also not necessarily a dependency: a document may
mention OpenSSL because of a license notice, as a tool used in guidance or
testing, or in a list of excluded components. Each mention therefore receives
an evidence level, as defined in Table~\ref{tab:evidence}. The results report
counts per level, so that a reader can see how strong the evidence behind each
figure is.

\begin{table}[htbp]
  \centering
  \caption{Evidence levels assigned to library mentions}
  \label{tab:evidence}
  \begin{tabular}{lp{0.72\textwidth}}
    \toprule
    Level & Meaning \\
    \midrule
    E1 & The library is explicitly described as a component of the product
      or module, and its version is stated. \\
    E2 & The library is explicitly described as a component, without a version. \\
    E3 & A dependency is implied but not stated, for example a derived module or a
      vendor-branded wrapper. \\
    E0 & The library is mentioned, but not as a dependency, for example in a
      license notice or as a test tool. Excluded from dependency counts. \\
    \bottomrule
  \end{tabular}
\end{table}

\section{AI-assisted document parsing}
\label{sec:ai}

Vendors may redact or generalize their descriptions of dependencies, so regular
expressions alone can miss implicit references and cannot judge the role of a
mention. A large language model is therefore used for two tasks: assigning an
evidence level and a normalized library and version to each candidate mention,
and searching a selected set of documents with no regular-expression match for
implicit references.

Model output is treated as a measured classifier rather than as ground truth:

\begin{enumerate}
  \item A random, stratified sample of mentions is annotated manually, and
    kept separate from the data used to develop the prompts.
  \item Precision and recall are reported on this sample, per evidence level and
    for regular expressions alone versus regular expressions with the model.
  \item Every model answer must quote the passage that supports it, and an
    answer whose quotation does not appear in the source document is rejected.
  \item The model identifier, prompts, and responses are stored, so that the
    analysis can be rerun from cached responses without calling the model.
\end{enumerate}

{Choose the model and decide the sample size and whether a second annotator
is available. Consider the cost of running the model on the full corpus.}

\section{Metrics}
\label{sec:metrics}

\paragraph{Prevalence (RQ1).} For each library and scheme, prevalence is the share of
certificates with a dependency at evidence level E1--E3, broken down by year,
vendor, product category, and, for CC, assurance level. Because a mention is not
evidence of evaluation, the analysis distinguishes the library's role where
documents allow it.

\paragraph{Certification lag (RQ2).} The lag is defined over certificates with a known
library version, which requires evidence level E1.

\begin{definition}[Certification lag]
  Let $c$ be a certificate whose product depends on library version $v$, let
  $d_{\mathrm{cert}}(c)$ be the date of the certification event, and let
  $d_{\mathrm{rel}}(v)$ be the upstream release date of $v$. The
  \emph{certification lag} of $c$ is
  \begin{equation}
    L(c) = d_{\mathrm{cert}}(c) - d_{\mathrm{rel}}(v).
    \label{eq:lag}
  \end{equation}
\end{definition}

The certification event is the initial certificate or a maintenance
or re-validation event, which are analyzed separately. Because calendar time does not
capture how far a product has fallen behind, the lag is complemented by the
\emph{release distance}, the number of upstream releases on the same branch
between $v$ and the newest release available at $d_{\mathrm{cert}}(c)$. As
Section~\ref{sec:libraries} shows, a validated provider version can differ from
the version of the library that loads it, so FIPS provider certificates are
reported separately.

{Decide how to treat vendor-patched versions, whose version strings can hide
backported fixes, and whether lag can also be computed for the forks, which may
have no comparable versioned releases.}

\paragraph{Business and supply-chain indicators (RQ1, RQ2).} These comprise library
adoption by vendor over time, the concentration of certificates on a few libraries
and vendors, the frequency of re-certification, migration between library families,
and, for FIPS, the state of the transition from FIPS~140-2 to 140-3.

\paragraph{Workflow quality (RQ3).} RQ3 is assessed by the extraction quality from
Section~\ref{sec:ai}, and by whether the complete analysis can be regenerated from
the pinned snapshot using the published notebooks.

\section{Analysis and visualization}
\label{sec:tooling}

The analysis is implemented in Python with \texttt{pandas}, and is
visualized with \texttt{matplotlib} and \texttt{seaborn}. Dependencies are
modeled as a bipartite graph of vendors and libraries using \texttt{networkx}, and
interactive views are built with \texttt{plotly}. The outputs are relationship
graphs, temporal lifecycle timelines, and summary tables.

\section{Reproducibility}
\label{sec:reproducibility}

The toolkit consists of Jupyter notebooks and processing scripts that run from a
pinned dataset snapshot with pinned dependencies. Each table and figure in
Chapter~\ref{ch:results} is generated by a notebook, so a reader can trace any
reported number to its source documents. Appendix~\ref{app:supplementary} describes
how to run it.
